%!TEX root = ../main.tex

\newpage

\begin{figure}[ht]
\centering
\begin{subfigure}[b]{0.90\textwidth}
    \centering
    \includegraphics[width=\textwidth]{figure1a.pdf}
\end{subfigure}
\vspace{0.5cm}
\begin{subfigure}[b]{0.90\textwidth}
    \centering
    \includegraphics[width=\textwidth]{figure1b.pdf}
\end{subfigure}
\caption{
\textbf{Microbial carbon utilisation prediction workflow.}
(A) Overview of data collection, genome annotation, and machine learning workflow.
We collected growth phenotypes (growth vs. no-growth) of microbial isolates across 242 carbon sources from four datasets (ATLeaf, Biolog, Marine, and Populus).
Genome sequences were annotated using RAST, KOFAM, and GapMind, and the resulting feature matrices were used to train the classifiers.
(B) Phenotypic distribution for the 15 shared carbon sources across the four datasets.
The bar plot displays sample counts and class distributions for each phenotype.
Solid bars indicate growth (positive) samples, and patterned bars indicate no-growth (negative) samples.
Class distributions varied across phenotypes and datasets.
}\label{fig:dataset_characteristics}
\end{figure}

\begin{figure}[ht]
\centering
\includegraphics[width=0.90\textwidth]{figure2.pdf}
\caption{
\textbf{Baseline mechanistic and phylogenetic predictions show limited accuracy.}
(A) Bar plot showing GapMind balanced accuracy across 15 phenotypes.
GapMind predictions were compared with experimental outcomes using both strict and permissive pathway completeness thresholds.
(B) Dot plot showing phylogeny-based nearest neighbour (k=3) balanced accuracy across 15 phenotypes under two evaluation scenarios: random split (top panel) and out-of-clade split (bottom panel).
Circles and boxes indicate null model performance (Identity and Bernoulli baselines) and triangles indicate nearest neighbour performance.
}\label{fig:baseline_predictions}
\end{figure}

\begin{figure}[ht]
\centering
\includegraphics[width=0.85\textwidth]{figure3.pdf}
\caption{
\textbf{ML models excel within datasets but struggle to generalise across datasets.}
(A) Box plot showing random cross-validation balanced accuracy for gradient-boosted tree classifiers (using KOFAM features) across 15 phenotypes.
The dotted line indicates the GapMind permissive baseline performance on the test set.
(B) Dot plot showing cross-dataset balanced accuracy when training on three datasets and testing on the fourth held-out dataset.
The dotted line indicates the GapMind baseline.
(C) Difference plot comparing model performance between regular cross-dataset splits and phylogenetically constrained in-clade cross-dataset splits.
}\label{fig:ml_performance}
\end{figure}

\begin{figure}[ht]
\centering
\includegraphics[width=\textwidth]{figure4.pdf}
\caption{
\textbf{Diagnosing GapMind misclassifications and ML generalisation failures.}
(A) Conceptual framework outlining causes of concordant samples (where GapMind matches experiment) and discordant samples (where GapMind contradicts experiment).
(B) Stacked bar plots showing per-phenotype (top) and per-dataset (bottom) misclassification patterns.
Each bar displays the proportions of true positives, true negatives, false positives (GapMind predicts complete pathway but no growth is observed), and false negatives (GapMind predicts incomplete pathway but growth is observed).
(C) Feature importance stability analysis using SHAP.
Top panel: number of stable features appearing in $\geq$70\% of model runs across random splits with 20 different random seeds.
Bottom panel: cross-dataset stability comparing individual dataset features to aggregated dataset (combination of the three other datasets) features.
Solid bars represent features common to the aggregated dataset, and patterned bars represent dataset-unique features.
}\label{fig:misclassification_diagnosis}
\end{figure}

\begin{figure}[ht]
\centering
\includegraphics[width=0.83\textwidth]{figure5.pdf}
\caption{
\textbf{Concordant-trained models learn mechanistic features and partially predict discordant samples.}
All models used the same architecture and feature representation as in Figure~3; only the training set was restricted to concordant samples.
(A) Balanced accuracy on held-out GapMind-concordant samples under random split (green line) and cross-dataset split (blue boxes) across 15 phenotypes.
(B) Cross-dataset feature stability analysis for concordant-trained models using SHAP (stable features are those that appear in $\geq$70\% of runs; see Methods).
Solid bars: features common to the aggregated dataset; patterned bars: dataset-unique features.
(C) Balanced accuracy of concordant-trained models applied to discordant samples.
Performance varied substantially across phenotypes; values near random indicate cases in which genome-derived features were insufficient to resolve the phenotype.
}\label{fig:concordant_analysis}
\end{figure}

\begin{figure}[ht]
\centering
\includegraphics[width=0.90\textwidth]{figure6.pdf}
\caption{
\textbf{Effects of data filtering and feature selection on model performance.}
(A) Identification of problematic samples across datasets.
Left: microbes showing zero growth on all substrates.
Middle: microbes showing universal growth despite incomplete GapMind pathways.
Right: Top 20 most misclassified genomes across all datasets.
(B) Grouped bar plot showing cross-dataset balanced accuracy after data filtering.
Blue: concordant-trained models.
Green: models after confidence-based filtering.
Red: models after removing bad genomes from panel A.
(C) Scatter plots of precision vs.\ recall.
Left: filtered data strategies.
Right: GapMind predictions alone (purple).
Each point represents average performance for one phenotype across cross-dataset splits.
(D) Effect of feature selection on model performance.
Right: precision-recall trade-off.
Right: precision-recall trade-off.
Circles: random splits; squares: cross-dataset splits.
}\label{fig:improvement_strategies}
\end{figure}

\begin{figure}[ht]
\centering
\includegraphics[width=0.95\textwidth]{figure7_gapmind_histidine_all_tests.pdf}
\caption{
\textbf{Sample size requirements for training generalisable models.}
Learning curves showing balanced accuracy as a function of training sample size for the histidine phenotype using GapMind-concordant data.
Columns represent evaluation scenarios: (left) random splits, (middle) cross-dataset splits, and (right) out-of-clade splits.
Rows represent test sample types: full (top), concordant (middle), and discordant (bottom).
Each point indicates performance at training sizes of 50, 100, 200, 500, or the full dataset. Performance on discordant test samples reflects cases in which genome-derived features are insufficient to fully encode the phenotype; lack of saturation with additional training data indicates representational, not statistical, limits.
}\label{fig:model_sample_size_requirements}
\end{figure}
